\documentclass{article}

\usepackage{amsmath}
\usepackage[margin=1in]{geometry}
\usepackage{graphicx}

\title{Phases Progress Summary}

\begin{document}

\section{Summary}
\subsection{Couplings}
Representative $(a_M, \varepsilon_r)$ points were chosen for each phase, with values
and couplings shown below (couplings in meV, $a_M$ in nm).
\begin{center}
\begin{tabular}{|c|c|c|c|c|c|c|c|c|c|}
    \hline
    Phase & $(a_M,\varepsilon_r)$ & $J_H$ & $J_{zz}^s$ & $J_{zz}^o$ & $J_+$ & $J_{+-}^s$ & $J_{++}^s$ & $J_{+-}^o$ & $J_{++}^o$\\
    \hline
    FM & $(9.0,5.0)$ & $-2.1$ & $-7.4$ & $-2.6$ & $-1.5 - 3.7i$ & $2.2 + 3.1i$ & $0.67$ & $0.22 + 0.33i$ & $0.82$\\
    \hline
    Stripe & $(6.0,5.0)$ & $-16.0$ & $-4.2$ & $-0.48$ & $4.2 - 26.0i$ & $1.8 + 20.0i$ & $-3.4$ & $0.11 + 1.0i$ & $0.44$\\
    \hline
    AFM-FE & $(10.0,11.0)$ & $22.0$ & $-8.0$ & $-2.2$ & $-14.0 + 12.0i$ & $2.0 - 8.8i$ & $11.0$ & $0.78 - 2.0i$ & $3.6$\\
    \hline
    AFM-AFE & $(7.0,11.0)$ & $14.0$ & $-36.0$ & $-8.5$ & $-11.0 - 14.0i$ & $1.3 + 10.0i$ & $15.0$ & $0.49 + 0.89i$ & $5.3$\\
    \hline
\end{tabular}
\end{center}
These couplings were then normalized by their L2 norm so that they would be of order 1, for better numerics,
so energies and temperatures will be plotted in the normalized units below.
\begin{center}
\begin{tabular}{|c|c|}
    \hline
    Phase & Norm (meV)\\
    \hline
    FM & 9.856\\
    \hline
    Stripe & 37.56\\
    \hline
    AFM-FE & 32.99\\
    \hline
    AFM-AFE & 47.32\\
    \hline
\end{tabular}
\end{center}

\subsection{Order Parameters}
As order parameters, the correlations of $\vec{S}$ and $\vec{\eta}$ were taken
at points in Fourier space, taking form:
\begin{align*}
    C_{\mathbf{k}}^{\alpha\beta} &\equiv \langle S_{-\mathbf{k}}^\alpha S_{\mathbf{k}}^\beta \rangle\\
    D_{\mathbf{k}}^{\alpha\beta} &\equiv \langle \eta_{-\mathbf{k}}^\alpha \eta_{\mathbf{k}}^\beta\rangle
\end{align*}
Due to $SU_s(2)$ symmetry, the spin correlation should be diagonal and isotropic, so only its trace
$C_{\mathbf{k}}\equiv C_{\mathbf{k}}^{\alpha\beta}\delta_{\alpha\beta}$ was used.

To decide on which $\mathbf{k}$ to measure at, the Fourier transform of the ground state was
taken, resulting in 1 or 2 peaks, and correlations were measured at those points, shown below
\begin{center}
\begin{tabular}{|c|c|c|}
    \hline
    Phase & Relevant $C_{\mathbf{k}}$ & Relevant $D_{\mathbf{k}}^{\alpha\beta}$\\
    \hline
    FM & $\mathbf{k}=\Gamma$ & $\mathbf{k}=\Gamma; \alpha,\beta=z$\\
    \hline
    Stripe & $\mathbf{k}=\pm M/2$ & $\mathbf{k}=M; \alpha,\beta=x,y$\\
    \hline
    AFM-FE & $\mathbf{k}=M$ & $\mathbf{k}=\Gamma; \alpha,\beta=x,y$\\
    \hline
    AFM-AFE & $\mathbf{k}=\pm 3K/4$ & $\mathbf{k}=M'; \alpha,\beta=x,y$\\
    \hline
\end{tabular}\\
\vspace{1ex}
\fbox{\includegraphics[width=0.3\textwidth]{plots/fbz.png}}
\end{center}
For phases where $\eta$ ordering was in-plane,
$D_{\mathbf{k}}^\parallel\equiv D_{\mathbf{k}}^{xx} + D_{\mathbf{k}}^{yy}$ was used
as an order parameter.

The Hamiltonian is $C_3$ symmetric, so $C_3$ rotated versions of the ground states above
will also form in the Monte Carlo simulation. To reduce simulation time, the system was
started with the system pinned to the ground states above with a strong bias field,
which was gradually turned off before measuring any observables. Since the system
was biased to one of the ground state orientations above, only the correlations for
that orientation were used.

To have an unbiased initialization, the system was also started in a random, high-$T$
state, and $T$ was gradually reduced to the target $T$ before measurements. In this case,
the system is not biased to a particular $C_3$ orientation of the ground state, so
the $C_3$ invariant correlations below were used:
\begin{align*}
    \tilde{C}_{\mathbf{k}} &\equiv C_{\mathbf{k}} + C_{C_3\mathbf{k}} + C_{C_3^2\mathbf{k}}\\
    \tilde{D}_{\mathbf{k}}^{zz} &\equiv D_{\mathbf{k}}^{zz} + D_{C_3\mathbf{k}}^{zz} + D_{C_3^2\mathbf{k}}^{zz}\\
    \tilde{D}_{\mathbf{k}}^{\parallel} &\equiv D_{\mathbf{k}}^{\parallel} + D_{C_3\mathbf{k}}^{\parallel} + D_{C_3^2\mathbf{k}}^{\parallel}
\end{align*}

\newpage
\section{Ferromagnet}
For the ferromagnet, temperature and energy are in units of meV, and order
parameters are the spin expectation (magnetization) and $\eta_z$ expectation.
\subsection{Zero Field}
\begin{itemize}
    \item State is ordered below $T=4.5$ meV ($\approx 52$ K), when initialized
    in ground state (high bias field), also see heat capacity spike at that temperature
    \item Transition temperature does not seem affected by system size
\end{itemize}
\begin{center}
    \frame{\includegraphics[width=0.85\textwidth]{plots/fm.png}}
\end{center}

\newpage
\section{Antiferromagnet-Ferroelectric}
For AFM-FE, the Hamiltonian is normalized, and the energy and T are in units
of the Hamiltonian (32.98 meV). The order parameters are the expectation of
norm of in-plane $\vec{\eta}$ and the Fourier spin correlation at M.
\subsection{Zero Field}
\begin{itemize}
    \item Spins don't show a clear transition, but in plane $\eta_{xy}$ does show a transition
    at $T=0.3\approx 114$ K
    \item Also see a peak in heat capacity at the same temperature
\end{itemize}
\begin{center}
    \frame{\includegraphics[width=0.85\textwidth]{plots/afm_fe.png}}
\end{center}

\newpage
\section{Stripe}
For stripe, the Hamiltonian is normalized, and the energy and T are in units
of the Hamiltonian (37.56 meV). The order parameters are the Fourier correlations
of in-plane $\vec{\eta}$ at M, and spins at M/2
\subsection{Zero Field}
\begin{itemize}
    \item See transition in spins, $\eta$, and heat capacity at T=0.5 ($\approx 436$ K)
\end{itemize}
\begin{center}
    \frame{\includegraphics[width=0.85\textwidth]{plots/stripe.png}}
\end{center}

\end{document}